\subsection{Calculus Vol2 - 7.1 Parametric Equations}

\url{https://openstax.org/books/calculus-volume-2/pages/7-1-parametric-equations}

\begin{enumerate}
  \item What is a parametric equation, and how does it differ from a rectangular equation $y=f(x)$?
  \item What is a parameter, and what does it typically represent in applications (e.g., time)?
  \item How are $x$ and $y$ expressed in terms of a parameter $t$ in parametric form?
  \item How does a point $(x(t), y(t))$ move as the parameter $t$ increases?
  \item Why are parametric equations useful for describing curves that are not functions?
  \item What does the orientation of a parametric curve mean, and how is it determined?
  \item Why can the same curve have multiple different parametric representations?
  \item How can you eliminate the parameter to convert parametric equations into a rectangular equation?
  \item What is the purpose of eliminating the parameter?
  \item How do restrictions on the parameter $t$ affect the portion of the curve that is traced?
  \item What are the parametric equations of a circle, and how do they relate to trigonometric functions?
  \item How can you verify that a parametric curve represents a known geometric shape?
  \item How do parametric equations provide more information than a rectangular equation (e.g., direction of motion)?
  \item What does it mean for a curve to be traced once, multiple times, or partially?
  \item How is the derivative $\frac{dy}{dx}$ computed for a parametric curve?
  \item Why is the formula $\frac{dy}{dx} = \frac{dy/dt}{dx/dt}$ valid?
  \item What happens when $\frac{dx}{dt} = 0$ in the formula for $\frac{dy}{dx}$?
  \item How can you find points where the tangent line is horizontal or vertical?
  \item What is a critical point for a curve defined parametrically?
  \item How can you find the equation of a tangent line to a parametric curve at a given value of $t$?
  \item How do you compute $\frac{dx}{dt}$ and $\frac{dy}{dt}$ for given parametric equations?
  \item What is the relationship between parametric equations and vector-valued functions?
  \item How can parametric equations describe curves that cannot be written as $y=f(x)$?
\end{enumerate}
